\clearpage\thispagestyle{empty}
\pdfbookmark[0]{Quellen zur Besprechung von Damoiseaus Mondtafeln}{damoiseau1824-sources}
\begin{center}\Large\bfseries Quellen zur Besprechung von Damoiseaus Mondtafeln\end{center}

Die Übertragung folgt Carl Friedrich Gauß’ Besprechung in den \emph{Werken}, Bd.~VI (Göttingen, 1874), S.~640--641. Der Erstabdruck erschien in den \emph{Göttingischen gelehrten Anzeigen}, Stück~23, 7.~Februar 1825, S.~225--227. Die \emph{Tables de la lune} des Barons de Damoiseau erschienen 1824 bei Bachelier in Paris.

KI-Übertragung aus dem historischen Deutsch in modernes Deutsch. KI-erstellter Satz.

\href{https://gdz.sub.uni-goettingen.de/id/PPN236007467}{Werkausgabe}. Erstdruck, S.~\href{https://images.p.sub.uni-goettingen.de/iiif/image/gdz:PPN385030444_1825_1:00000225/full/1650,/0/default.jpg}{225}, \href{https://images.p.sub.uni-goettingen.de/iiif/image/gdz:PPN385030444_1825_1:00000226/full/1650,/0/default.jpg}{226}, \href{https://images.p.sub.uni-goettingen.de/iiif/image/gdz:PPN385030444_1825_1:00000227/full/1650,/0/default.jpg}{227}. \href{https://books.google.it/books?id=uixOAAAAYAAJ&pg=PR3}{Damoiseaus Mondtafeln}.

\clearpage
